%======================================================================
\section{Examples}
\label{sec:examples}

All real-market examples in the paper are fitted from one frozen
snapshot: SPY and NVDA listed options, eight expiries each, from
\MacNvdaOneDayDays\ calendar days to \MacNvdaLongTYears\ years, captured
after the close of the 2026-08-03 US session from a delayed consolidated
feed, with spots $\MacSpySpot$ (SPY) and $\MacNvdaSpot$ (NVDA).  Fits
use the haircut target \eqref{eq:haircut}, the logistic chart, the
shipped order $N=16$ under the guard \eqref{eq:orderguard}, and the
calendar rows \eqref{eq:calrows}; nothing was tuned per example.  One
consequence of the recipe should be recalled from
\cref{sec:calibration} before reading the figures: SPY spreads are all
tighter than the haircut width, so every SPY band degenerates to a mid
target and the SPY panels draw raw bid--ask quotes, while the NVDA
nodes carry live bands.  Across all \MacSnapNodeCount\ fitted nodes the
median quote-fit error is \MacSnapMedianRmsBp\ vol bp and the worst
\MacSnapWorstRmsBp\ vol bp (the shortest NVDA node, discussed below).
Provenance, quality counts, and regeneration instructions are in
\cref{app:data}; the synthetic examples of this section are exactly
specified there as well.

\subsection{The SPY gallery}

\Cref{fig:spy_gallery} shows all eight SPY expiries: quotes as raw
bid--ask whiskers with their mids, fitted slices through them,
per-expiry error annotated.  The gallery's purpose is not to impress
with residuals --- a flexible basis fitting tight liquid mids
\emph{should} produce small residuals --- but to exhibit the shape
vocabulary of one underlier across maturities:
the steep short-dated put skew relaxing into the long-dated smile, the
at-the-money level drifting with the term structure, the wings beyond
the outer quotes following the two-parameter exponential-tail prior of
\cref{sec:tails}.  Median error across the eight slices is
\MacSpyGalleryMedianRmsBp\ vol bp, worst \MacSpyGalleryWorstRmsBp.

\begin{figure}[htbp]
\centering
\paperfig{\textwidth}{fig_spy_gallery.pdf}
\caption{Eight SPY expiries from the frozen snapshot, fitted with one
recipe.  Each panel: raw bid--ask quotes (whiskers) and mids, fitted
LQD slice (line), per-expiry rms in vol bp (annotation).  The haircut
bands do not appear because none survive on SPY --- every spread is
tighter than $2h$, so the fit target is the mid everywhere
(\cref{sec:calibration}).  Quoted strike coverage shrinks and steepens
toward the short end; everything beyond the outermost quotes in each
panel is the model's exponential-tail extrapolation, labeled as prior,
not data, per \cref{sec:tails}.}
\label{fig:spy_gallery}
\end{figure}

\subsection{One node under the microscope}

\Cref{fig:spy_node} takes the December 2026 SPY expiry --- the slice
already used in \cref{fig:lqdchart,fig:butterfly,fig:jacobian} --- and
reads it three ways: fit against the raw bid--ask quotes, residual
ladder in vol bp per strike, and the fitted density against an
ATM-matched normal.  The node fits \MacSpyDecNQuotes\ quotes to
\MacSpyDecRmsBp\ vol bp rms.  Its exact handles (\cref{prop:atm}): ATM
volatility \MacSpyDecAtmVolPct\%, skew $\sigma'(0)=\MacSpyDecSkew$,
curvature $\sigma''(0)=\MacSpyDecCurv$.  The fitted forward rank is
$u_*=\MacSpyDecForwardRankPct$ --- below one half, and, what actually
carries the skew, below the matched flat-Black percentile $\PhiN(d)$:
this is the digital mismatch $\Delta<0$ of \cref{prop:atm} in live
numbers, one sign with the measured $\sigma'(0)<0$, and the contrast
with the symmetric slice of \cref{sec:ledger} (whose forward rank sits
\emph{above} one half with $\Delta\approx0$) shows which effect owns
which comparison: the martingale shift moves $u_*$ against one half,
the skew moves it against $\PhiN(d)$, and the second comparison is the
one the smile trades on.  The var-swap strike \eqref{eq:varswap} is
\MacSpyDecVarSwapPct\%, well above the ATM volatility --- as expected
for a slice with this much left tail, since the log contract weights
exactly the wings the smile lifts.  The density panel shows where that
premium lives: mass moved from the right shoulder into the left tail
relative to the matched normal.

\begin{figure}[htbp]
\centering
\paperfig{\textwidth}{fig_spy_node.pdf}
\caption{The SPY December 2026 node.  Panel A: raw bid--ask quotes and
the fitted slice (no haircut bands survive on SPY ---
\cref{sec:calibration}).  Panel B: the residual ladder --- signed fit
error versus mid in vol bp at each quote, drawn inside the bid--ask
envelope; rms \MacSpyDecRmsBp\ vol bp.  Panel C: the fitted density
against a normal matched at the money: the skew is a transfer of mass
into the left tail, paid for by a thinner right shoulder --- the same
fact the negative digital mismatch $\Delta<0$ of \cref{prop:atm} states
in one number.}
\label{fig:spy_node}
\end{figure}

\paragraph{A ticket priced by hand.}
The ledger makes a price auditable with a pencil, which is worth doing
once in public, on a contract that actually trades.  On this node,
price the listed December \MacTicketStrike\ call: log-moneyness
$k=\MacTicketK$ against the node's forward.  The strike root gives rank
$u_k=\MacTicketRankPct$: the option finishes in the money with
probability $1-u_k$.  The two ledger entries \eqref{eq:call} read
$G(z_k)=\MacTicketShare$ --- the share of the forward carried by
outcomes above the strike --- and
$e^{k}(1-u_k)=\MacTicketCash$ --- the strike bill paid on that event ---
so
\[
c(\MacTicketK)=\MacTicketShare-\MacTicketCash=\MacTicketPrice,
\]
or \MacTicketDollarPrice\ per share in dollars through the node's
forward and discount, for an implied volatility of \MacTicketIvPct\%.
Reproducing these five numbers from the snapshot file is a one-screen
exercise and the fastest full-stack audit of an implementation the
author knows.

\subsection{NVDA, including the shortest node}

Single-name smiles are heavier-tailed and coarser-quoted --- and, per
\cref{sec:calibration}, NVDA is where the haircut band mechanics
actually engage, with roughly half the quotes carrying live bands on
the featured nodes.  \Cref{fig:nvda_nodes} shows two NVDA expiries: the
shortest node (\MacNvdaOneDayDays\ calendar days from the reference
date) and the longest, December 2027.
The long node behaves like SPY's slices with wider tails
(\cref{fig:tails}): \MacNvdaLongNQuotes\ quotes, rms
\MacNvdaLongRmsBp\ vol bp, fitted at full order.  The shortest node is
the order guard \eqref{eq:orderguard} working in public:
\MacNvdaOneDayNQuotes\ retained quotes admit effective order
$N_{\mathrm{eff}}=\MacNvdaOneDayEffOrder$, and the fit lands at
\MacNvdaOneDayRmsBp\ vol bp rms --- the largest error in the snapshot,
on the thinnest, jumpiest strip, with the fewest parameters: exactly
where a fitter should be least confident, and visibly so in the
figure's wider bands.

\begin{figure}[htbp]
\centering
\paperfig{\textwidth}{fig_nvda_nodes.pdf}
\caption{Two NVDA nodes from the same recipe, with the haircut bands
live (unlike SPY's, NVDA's spreads are wide enough to survive the
shrink \eqref{eq:haircut}).  Left: the shortest expiry
(\MacNvdaOneDayDays\ calendar days) --- \MacNvdaOneDayNQuotes\ quotes,
order-guarded to $N_{\mathrm{eff}}=\MacNvdaOneDayEffOrder$, rms
\MacNvdaOneDayRmsBp\ vol bp; the guard trades inter-quote flexibility
for stability on a strip that cannot support sixteen modes
(\cref{sec:orderguard}).  Right: the December 2027 expiry at full order
--- \MacNvdaLongNQuotes\ quotes, rms \MacNvdaLongRmsBp\ vol bp --- with
the single-name tail scales visibly heavier than SPY's
(\cref{fig:tails}).}
\label{fig:nvda_nodes}
\end{figure}

\subsection{Order control: the synthetic double-hump}

What does basis order buy?  \Cref{sec:validity} said the family can
approximate bimodal laws; this example measures it.  The target is a
known mixture of two lognormal components (parameters in
\cref{app:data}), martingale-normalized, quoted as noiseless implied
volatilities on a realistic strike grid --- an approximation benchmark,
not a market fit.  At $N=16$ the fit reproduces the target smile to
\MacHumpFitRmsBp\ vol bp rms (maximum error \MacDhMaxErrBpHigh\ vol bp)
and recovers the bimodal density \emph{faithfully}: modes at
\MacDhModeLeft\ and \MacDhModeRight\ against the true
\MacDhTrueModeLeft\ and \MacDhTrueModeRight, trough and asymmetry
intact (\cref{fig:doublehump}).  At $N=6$ (\cref{fig:order_control})
the smile moves little --- the two fits agree in implied volatility to
\MacHumpIvAgreeBp\ vol bp on a dense grid, with the larger gaps sitting
out in the wings beyond the quoted ladder --- and the density keeps
both modes (\MacDhModeCountLow, at this mixture separation) but is
\emph{distorted}: the peaks overshoot, the trough is half-filled, the
mode locations drift.  The honest scalar carrier of the difference is
the $L^1$ distance to the true density, which more than doubles:
\MacHumpLOneNSixteen\ at $N=16$ against \MacHumpLOneNSix\ at $N=6$.

The pair of figures is the paper's cleanest exhibit on information
versus convention.  The \emph{quotes} barely distinguish the two fits;
the \emph{densities} differ measurably --- not, at this separation, by
the existence of the modes, but by their fidelity: shape, depth, and
placement of the structure between the quotes.  Below the quote
spacing, that fidelity is chosen by the basis order, not by the market
--- so a desk that trades density-sensitive structures (digitals,
tight butterflies) around events should set $N$ as a conscious
convention, and a desk that only reads smiles may never notice the
difference.  Order is power to assert structure between quotes; the
guard of \cref{sec:orderguard} caps that power when quotes are few.

\begin{figure}[htbp]
\centering
\paperfig{\textwidth}{fig_doublehump.pdf}
\caption{The double-hump benchmark at $N=16$.  Panel A: target implied
volatilities from the bimodal mixture and the fitted slice ---
agreement to \MacHumpFitRmsBp\ vol bp rms; note the smile shows
bimodality only as gentle shoulder structure.  Panel B: the true
mixture density (shaded) and the fitted LQD density (dashed): both
modes, the trough, and the asymmetry are recovered, with fitted mode
locations \MacDhModeLeft\ / \MacDhModeRight\ against the true
\MacDhTrueModeLeft\ / \MacDhTrueModeRight; positivity was never at risk
while the high modes carved the trough.}
\label{fig:doublehump}
\end{figure}

\begin{figure}[htbp]
\centering
\paperfig{\textwidth}{fig_order_control.pdf}
\caption{The same target at $N=6$.  Panel A: the $N=6$ and $N=16$
smiles nearly coincide (maximum gap \MacHumpIvAgreeBp\ vol bp, the
larger gaps beyond the quoted ladder) --- the quotes can barely tell
the fits apart.  Panel B: the densities can: the $N=6$ fit keeps both
modes but overshoots the peaks, half-fills the trough, and drifts the
mode locations --- $L^1$ error to the truth \MacHumpLOneNSix\ against
\MacHumpLOneNSixteen\ at $N=16$, more than double.  Between the quotes,
order is a convention, and this is what choosing it asserts.}
\label{fig:order_control}
\end{figure}

\subsection{The exact audit}

The fifth example is the oldest kind: against Euler.
\Cref{fig:exact} (\cref{sec:constspeed}) already confronted the
implementation with the constant-speed slice's closed forms across the
admissible scale range --- transport map to \MacExactMapWorst, shift to
\MacExactMuWorst\ (growing lawfully toward the wall), martingale check
to \MacExactMartWorst.  The example is listed again here because of
what it covers: quadrature, tail corrections, normalization,
interpolation, and root-solving in one line of algebra with no free
parameters.  A wrong grid constant, a mis-anchored cumulative sum, or a
stale cache moves any of these agreements by many orders of magnitude;
it is the cheapest regression test the model owns, and the recommended
first test of any reimplementation from this paper.
